\scanpage{41}
One can also introduce the groupoid
$C_{1,1}^{\mathrm{cat}\,\prime}$, obtained from the \inserted{groupoid}
$C_{1,1}^{\mathrm{ext}}$ by \unsure{passing to the|taking a}
quotient by the universal automorphism $-\mathrm{id}$
--- corresponding to the center \inserted{$\{\pm1\}$} of the
\unsure{fundamental|fundamental} groups of this groupoid\unsure{.} --- one obtains
a connected groupoid equivalent canonical\textsuperscript{ly}
to that of torsors under \inserted{$\mathrm{Gl}(2,\mathbb Z)'=$}
$\mathrm{Gl}(2,\mathbb Z)/(\pm1)$,
\unsure{from which|whence} one now obtains an
\uline{equivalence} of \uline{categories}
\[
(55)\qquad C_{1,1}^{\mathrm{ext}\,\prime}\xrightarrow{\approx}C_{0,3}
\]
making commutative

\begin{samepage}
(56)\quad
\diagram{\begin{tikzcd}[ampersand replacement=\&,column sep=3.8em,row sep=3.4em]
C_{1,\unsure{0|1}}^{\mathrm{ext}\,\prime}\arrow[r,dash,"\approx"]\arrow[d] \& C_{0,3}\arrow[d]\\
\operatorname{Tors}\bigl(\mathrm{Gl}(2,\mathbb Z)'\bigr)\arrow[r,dash,"\approx"] \&
\operatorname{Tors}\bigl(\mathrm{GP}(1,\mathbb Z)\bigr)
\end{tikzcd}}
\end{samepage},

where $\operatorname{Tors}(\mathrm{Gl}(2,\mathbb Z)')\xrightarrow{\approx}
\operatorname{Tors}(\mathrm{GP}(1,\mathbb Z))$ is
the tautological functor associated to the isomorphism
\[
(57)\qquad \mathrm{Gl}(2,\mathbb Z)'\xrightarrow{\approx}\mathrm{GP}(1,\mathbb Z).
\]

\scanpage{42}
\mspage{309}
In short, I am in the process of making a
topological isotopic paraphrase, in terms
of the fundamental groups with loops, of
\uline{Legendre}'s theory of the \unsure{relations} between
elliptic curves and rational curves.
I would like now to make explicit the
\uline{Legendre} equivalence (55), which had
been obtained thanks to the knowledge of the
appropriate Teichmüller groups relative
to the two sides and of an ad-hoc isomorphism
between these two groups,
directly in terms of the
fundamental groups of \unsure{type|types} $\Pi_{0,3}$ \unsure{or|and} $\Pi_{1,1}$.

If $\pi'$ is a group with loops of type
\unsure{$1.1$|$1,1$}, it is appropriate first to note that
there is in $\operatorname{Autext}_{\mathrm{lac}}(\pi')$ a canonical element
of order $2$ --- in fact the unique nontrivial element
of the center --- denoted as is only right $(-1)_{\pi'}$.
\margin{\uline{N.B.} \unsure{$(-1)_{\pi'}$|$\operatorname{Aut}_{\pi'}$} preserves the orientation i.e.
has \unsure{multiplier|multiplicity} $+1$}

\scanpage{43}
Thus, one can consider $\pi'$ as endowed \inserted{canonical\textsuperscript{ly}}
with an \deleted{\unsure{action|other}} outer operation of
$G=\mathbb Z/2\mathbb Z$ on it, whence a
canonical extension
\[
(58)\qquad 1\longrightarrow\pi'\mathrel{\text{---}}\mathcal E\longrightarrow\mathbb Z/2\mathbb Z\longrightarrow1.
\]

The topological model teaches us
that there are exactly three $\pi'$-
conjugacy classes of splittings of this
extension, corresponding to the three
points of order $2$ of a topological group
$E\simeq\mathbf U\times\mathbf U$ --- the three fixed points of
$x\mapsto-x$ in $E\setminus\{0\}$. The theory of
``\unsure{drilling} holes \unsure{simultaneously}'', relative to a
\unsure{situation} of a \inserted{finite} group of outer operators
\unsure{of a|of a} group with loops (theory which
\unsure{even} in the discrete case has not really
been \unsure{written} --- for the moment there are
conjectural \unsure{statements}) then makes it possible
\margin{but here, one is in a case \illegible{2} can get by with
drillings \illegible{1} of one hole at a time, \unsure{repeated} successively\unsure{...}}

\scanpage{44}
\mspage{310}

to construct canonically an outer group \inserted{with loops $\widetilde{\pi}'$} of type $\Pi_{1,4}$, with an outer operation of $G=\mathbb Z/2\mathbb Z$ on this group, and a subset $J$ \inserted{$=I\setminus\{i_0\}$} of cardinality $3$ of $I=I(\widetilde{\pi}')$ (which is of cardinality $4$), $J$ \inserted{stable under $G$ i.e. \unsure{$i_0$ fixed by $G$}}, and an isomorphism \unsure{of the|of} group $\widetilde{\pi}'(J)$ \inserted{of type $\Pi_{1,1}$} obtained by ``plugging the holes at $J$'' with $\pi'$
\[
(59)\qquad \widetilde{\pi}'(J)\simeq\pi',
\]
isomorphism which is compatible with the operations of $G$. In the topological model, $\widetilde{\pi}'$ is the categorical fundamental group of \inserted{the space} $E\setminus {}_2E$, \inserted{space} on which $G$ operates freely.

\margin{in fact, $G$ operates \unsure{trivially} on $I=I(\widetilde{\pi}')$}

The extension $\widetilde{\mathcal E}$ of $G$ by $\widetilde{\pi}'$ \unsure{corresponding|corresponding} to the action of \deleted{$\unsure{\widetilde{\mathcal E}}$} $G$ on $\widetilde{\pi}'$
\begin{samepage}
\noindent (60)\quad
\diagram{\begin{tikzcd}[ampersand replacement=\&]
1 \arrow[r, dash] \& \widetilde{\pi}' \arrow[r, dash] \& \widetilde{\mathcal E} \arrow[r, dash] \& \mathbb Z/2\mathbb Z \arrow[r] \& 1
\end{tikzcd}}
\end{samepage},
\inserted{which} is identified \unsure{thus} with the mixed fundamental group of $(E\setminus {}_2E,G)$, is thus identified \unsure{also}-

\scanpage{45}

with the fundamental group of
\[
(61)\qquad V=(E\setminus {}_2E)/\{\pm1\}
\]
which is a topological surface of type $(0,4)$.

In other words, $\widetilde{\mathcal E}$ can be considered as a group with loops of type $(0,4)$, for a loop structure \unsure{easy} to describe in terms of that of $\widetilde{\pi}'$, and for which one has a canonical isomorphism
\[
(62)\qquad I(\widetilde{\mathcal E})\simeq I(\widetilde{\pi}')/G\simeq I(\widetilde{\pi}').
\]
Thus one has an $i_0\in I(\widetilde{\mathcal E})$, which allows one to ``plug the hole at $i_0$'' to obtain an \inserted{actual} group with loops of type $(0,3)$, i.e. \deleted{$\unsure{\pi=}$}
\[
(63)\qquad \pi=\mathcal L(\pi')
\]
Taking up in the reverse direction the relation between $\pi$ and $\pi'$, one can proceed as follows: to an \uline{actual} group \inserted{with loops} of type $\unsure{(0,3)}$, one associates by

\scanpage{46}
\mspage{311}

the ``drilling holes'' functor an \uline{outer} group with loops \inserted{$\widetilde{\pi}(\unsure{=\widetilde{\mathcal E}|\in\widetilde{\mathcal E}})$} of type $(0,4)$, endowed with an element $i_0\in I(\widetilde{\pi})$ such that \deleted{\illegible{1}}
\[
(64)\qquad I(\widetilde{\pi})\simeq\{i_0\}\amalg\underbrace{I(\pi)}_{\text{cardinality }3}
\]
One then has ($T$ being the orientation module of $\pi$)
\[
(65)\qquad \widetilde{\pi}_{\mathrm{ab}}\simeq T^{I(\widetilde{\pi})}\qquad (\widetilde{\pi}_{\mathrm{ab}})_2\simeq(\mathbb Z/2\mathbb Z)^{I(\widetilde{\pi})}
\]
and one has an (outer) homomorphism
\[
(66)\qquad \unsure{\widetilde{\pi}}\longrightarrow\mathbb Z/2\mathbb Z
\]
canonical, defined by the \unsure{condition} of being ``\unsure{ramified}'' at all the $i\in I(\widetilde{\pi})$ --- i.e. corresponding to the homomorphism
\[
(67)\qquad (\widetilde{\pi}_{\mathrm{ab}})_2\xrightarrow{\sim}(\mathbb Z/2\mathbb Z)^{I(\widetilde{\pi})}\longrightarrow\mathbb Z/2\mathbb Z
\]
\uline{sum} of the coefficients. The kernel $\widetilde{\pi}'$ of (66) is a group with loops of type $(1,4)$ --- but it is not \unsure{quite} an outer group --- it is not defined ``\unsure{modulo|up to} inner \unsure{automorphisms}'', \deleted{but} only modulo (in addition) the outer operation of \unsure{symmetry} of $\mathbb Z/2\mathbb Z$


\scanpage{47}

on it. If one \unsure{chooses} the group $\widetilde{\pi}$ (in its ``outer class'') -- one will in any case have an exact sequence
\[
(68)\qquad 1\longrightarrow\widetilde{\pi}'\longrightarrow\widetilde{\pi}\mathrel{\text{---}}\mathbb Z/2\mathbb Z\longrightarrow 1
\]
and $\widetilde{\pi}'$ is a group with loops, with
\[
(69)\qquad I(\widetilde{\pi}')\simeq I(\widetilde{\pi})
\]

By plugging the holes corresponding to $I(\pi)\subset I(\widetilde{\pi})\unsure{\simeq|=}I(\widetilde{\pi}')$, which is of cardinality $3$, one obtains a \deleted{outer} group of type $(1,1)$, with operation of $\mathbb Z/2\mathbb Z$ on it \inserted{which is none \unsure{other} than the canonical operation}. But depending on the indeterminacies of the choices made, it is a little less than an outer group of type $(1,1)$ i.e. an object of $C^{\mathrm{ext}}_{1,1}$ -- it is such an \unsure{element}, ``\uline{modulo} symmetry'' -- i.e. an object of $C^{\mathrm{ext}\unsure{\prime}}_{1,1}$. To have an object of $C^{\mathrm{ext}}_{1,1}$, one must have ``chosen a kernel'' $\widetilde{\pi}'$ for \unsure{the homomorphism} which is outer $\widetilde{\pi}\unsure{\longrightarrow|\mathrel{\text{---}}}\mathbb Z/2\mathbb Z$, as an outer group -- this is the choice of an object

\scanpage{48}
\mspage{312}

in a $(\mathbb Z/2\mathbb Z)$-gerbe, i.e. \inserted{\unsure{in}} a connected groupoid admitting $\mathbb Z/2\mathbb Z$ as fundamental group.
\margin{\unsure{3}}

\unsure{Modulo} a more \unsure{careful} write-up, it seems to me that these conceptual reflections do indeed prove and without \unsure{a flaw|flaws} that one has a canonical isomorphism (at least \unsure{modulo} \unsure{inner|inner} automorphisms)
\[
(70)\qquad \mathcal T_{1,1}/\{\pm1\}\xrightarrow{\sim}\mathfrak S_{0,3}
\]
\margin{in fact, one can make \unsure{it|it} precise \unsure{modulo inner automorphisms} arising \unsure{from|from the} $\mathfrak S^{\unsure{1+|i,+}}_{0,3}\simeq\Pi_{0,3}$}
-- the indeterminacy \unsure{modulo inner automorphisms} arising \unsure{quite} certainly from the choice of an isomorphism between $\mathcal L(\Pi_{1,1})$ and $\Pi_{0,3}$.

This fact is quite independent, \deleted{of} it seems to me, of the fact that \unsure{the homomorphism} which is natural
\[
\mathcal T_{1,1}\xrightarrow{\unsure{\sim}}\mathrm{Gl}(2,\mathbb Z)
\]
is an isomorphism. I see no objection to identifying $\mathcal T_{1,1}$ purely and simply with $\mathrm{Gl}(\unsure{2|1},\mathbb Z)$ on occasion\unsure{,|;} on the other hand it is more questionable to identify $\mathfrak S_{0,3}$ with $\mathcal T_{1,1}/\pm1$ \unsure{or with} $\mathrm{Gl}(2,\mathbb Z)'$

\scanpage{49}

except only as an outer group, because of the indeterminacy by inner automorphisms in (70).

From the point of view of analytic modular multiplicities, one has a canonical \unsure{morphism}
\[
(71)\qquad \mathbb M^{\mathrm{an}}_{1,1}\longrightarrow
\underbrace{\unsure{\mathbb U|\mathbb M}^{\mathrm{an}}_{0,\unsure{3}}}_{\inserted{\substack{\text{\unsure{complex|complete} analytic curve}\\\text{of type $(0,3)$ ``universal''}}}}
\simeq \mathbb M^{\prime\mathrm{an}}_{0,4}\simeq(U^{\mathrm{an}}_{0,3},\mathfrak S_3)
\]
\margin{with a little more care, one could include in the \unsure{reflections} also complex conjugation, \unsure{so as|in order} not to have \unsure{in|in} (73) to restrict oneself to the subgroups $\mathcal T^+$ and $\mathfrak S^+$ \unsure{in question}}
\deleted{whose effect on} which \unsure{induces} an isomorphism
\[
(72)\qquad \mathbb M^{\mathrm{an}\prime}_{1,1}\xrightarrow{\sim} M^{\mathrm{an}}_{0,3},
\]
and the effect of (71), on the outer fundamental groups is none other than the \unsure{categorical|outer} homomorphism
\[
(73)\qquad \mathcal T^+_{1,1}\longrightarrow\mathfrak S^+_{0,3}
\]
induced by (70) -- is it possible to make it precise in a useful way, as an actual \unsure{homomorphism}, not merely outer? There is a canonical universal \unsure{covering|covering} $\widetilde{\mathbb M}^{\mathrm{an}}_{1,1}$ of $\mathbb M^{\mathrm{an}}_{1,1}$, provided by the space \unsure{of|of the}

\scanpage{50}
\mspage{313}

Teichmüller (the \unsure{Poincaré} half-plane $D$ of \S\ 38) -- relative to this universal covering one does indeed have a canonical isomorphism
\[
(74)\qquad \pi_1(\unsure{M|\mathbb M}^{\mathrm{an}}_{1,1})\xrightarrow{\sim}\mathrm{Sl}(2,\mathbb Z)
\]
(not merely outer), which is moreover none other, essentially, than (50) -- (restricted to the subgroup of index $2$ $\mathcal T^+_{1,1}$ of $\mathcal T_{1,1}$). One does indeed have a canonical morphism.
\[
(75)\qquad D\longrightarrow U^{\mathrm{an}}_{0,3}\qquad(\simeq \unsure{M|\mathbb M}_{1,1}[2]')
\]
\margin{For the definition in \unsure{form|terms} of \unsure{this morphism}, cf. \S\quad}
compatible with a homomorphism on the groups of operators
\[
(76)\qquad \mathrm{Sl}(2,\mathbb Z)'\longrightarrow
\underbrace{\mathrm{Sl}(2,\mathbb Z/2\mathbb Z)\xrightarrow{\sim}\mathfrak S_3}_{\inserted{\substack{\text{\unsure{this isomorphism|the isomorphism} \unsure{which} is appropriate}\\\text{here is made explicit in}\\\text{\S\quad, page\quad}}}}
\]
so that $D$ appears as a universal covering of $U^{\mathrm{an}}_{\unsure{0,3|3}}$, hence also of $(U^{\mathrm{an}}_{0,3},\mathfrak S_3)$, and $\mathrm{Sl}(2,\mathbb Z)'$ as the mixed $\pi_1$ relative to this particular universal covering. To \inserted{\unsure{be able}} to identify it with $\pi_1((U^{\mathrm{an}}_{0,3},\mathfrak S_3);P=\unsure{-\bar{\jmath}|-j})$, it
